% EXPERIMENTAL METHODOLOGY
% Last generated: 2026-01-15T05:34:34.872Z

\subsection{Experimental Design}

We employ Monte Carlo simulation with systematic parameter sweeps to investigate our research questions. Each experiment configuration is run multiple times with different random seeds to establish statistical distributions.

\subsubsection{Parameter Sweeps}

For each research question, we vary a single parameter while holding others constant (ceteris paribus):

\begin{table}[htbp]
\centering
\caption{Experimental parameter configurations}
\label{tab:experiments}
\small
\begin{tabular}{@{}llll@{}}
\toprule
Experiment & Parameter & Values & Runs \\
\midrule
% AUTO-UPDATED: experiment list
Quorum Sensitivity & \texttt{quorumPercent} & 1--20 & 50 \\
DAO Scale & \texttt{totalMembers} & 25--500 & 50 \\
Voting Systems & \texttt{votingSystem.type} & majority, quad, conviction & 50 \\
Participation & \texttt{votingActivity} & 0.1--0.8 & 50 \\
\bottomrule
\end{tabular}
\end{table}

\subsubsection{Baseline Configuration}

All experiments use a common baseline derived from Compound governance:

\begin{itemize}
    \item Members: 100 (unless swept)
    \item Token distribution: Power-law ($\alpha = 1.5$)
    \item Quorum: 4\% (unless swept)
    \item Voting period: 7 days
    \item Timelock: 2 days
    \item Proposal frequency: 0.5/day
    \item Simulation length: 500 steps (days)
\end{itemize}

\subsection{Metrics}

We capture the following metrics for each simulation run:

\subsubsection{Primary Metrics}

\begin{description}
    \item[Proposal Pass Rate] Fraction of proposals that achieved quorum and majority support
    \item[Average Turnout] Mean participation rate across all votes
    \item[Final Gini] Gini coefficient of token holdings at simulation end
    \item[Treasury Efficiency] Funds deployed to successful projects vs. total treasury
\end{description}

\subsubsection{Secondary Metrics}

\begin{description}
    \item[Quorum Failure Rate] Proposals failing specifically due to quorum
    \item[Delegation Concentration] Herfindahl-Hirschman Index of delegated power
    \item[Time to Decision] Average steps from proposal to resolution
\end{description}

\subsection{Statistical Analysis}

\subsubsection{Summary Statistics}

For each parameter value, we compute:
\begin{itemize}
    \item Mean, median, standard deviation
    \item 95\% confidence intervals
    \item Interquartile range
\end{itemize}

\subsubsection{Hypothesis Testing}

We employ:
\begin{itemize}
    \item \textbf{ANOVA}: For comparing means across categorical parameters
    \item \textbf{Regression}: For continuous parameter relationships
    \item \textbf{Kruskal-Wallis}: For non-normal distributions
\end{itemize}

Significance level $\alpha = 0.05$ with Bonferroni correction for multiple comparisons.

\subsubsection{Effect Size}

We report Cohen's $d$ for pairwise comparisons:
\begin{equation}
d = \frac{\bar{x}_1 - \bar{x}_2}{s_{\text{pooled}}}
\end{equation}

\subsection{Validation}

\subsubsection{Internal Validity}

We ensure internal validity through:
\begin{itemize}
    \item Deterministic seeding for reproducibility
    \item Multiple runs per configuration (minimum 30)
    \item Systematic parameter variation
\end{itemize}

\subsubsection{External Validity}

Limitations on external validity:
\begin{itemize}
    \item Agent models are simplified representations
    \item No adversarial behavior (MEV, governance attacks)
    \item Static preferences (no learning)
\end{itemize}

We discuss these limitations in Section~\ref{sec:limitations}.

\subsection{Computational Resources}

% AUTO-UPDATED: computation stats
Experiments were run on:
\begin{itemize}
    \item Hardware: Consumer workstation (details in Appendix)
    \item Concurrency: 4-8 parallel simulations
    \item Total compute time: See reproducibility manifest
    \item Total simulation runs: 6
\end{itemize}
